\documentclass[../../main.tex]{subfiles}

\begin{document}

\section{连通的拓扑空间}\label{section:ch11-连通}

拓扑空间 $X$ 的两个非空开集称为分离 $X$, 若它们不交且它们的并是 $X$.

若一个拓扑空间不能被这样的一对子集分离, 则称为是连通的. 由于开集的补是闭的, 一个空间的分离中的每个开集也是闭的. 因此一个拓扑空间是连通的当且仅当仅有的既开又闭的子集是全空间和空集.

$X$ 的子集 $E$ 称为连通的, 若它是连通的拓扑子空间. 因此 $X$ 的子集 $E$ 是连通的当且仅当不存在 $X$ 的开集 $\mathcal{O}_{1}$ 与 $\mathcal{O}_{2}$ 使得
\begin{align*}
\mathcal{O}_{1}\cap E\neq\varnothing,\quad\mathcal{O}_{2}\cap E\neq\varnothing,\quad E\subseteq\mathcal{O}_{1}\cup\mathcal{O}_{2}\text{ 且 }E\cap\mathcal{O}_{1}\cap\mathcal{O}_{2}=\varnothing.
\end{align*}

\begin{proposition}\label{proposition:ch11-连通连续象}
令 $f$ 为从连通空间 $X$ 到拓扑空间 $Y$ 的连续映射. 则它的象 $f(X)$ 是连通的.
\end{proposition}

\begin{proof}
观察到 $f$ 是 $X$ 映上拓扑空间 $f(X)$ 的连续映射, 其中 $f(X)$ 具有承袭自 $Y$ 的子空间拓扑. 我们用反证法. 假设 $f(X)$ 不是连通的. 令 $\mathcal{O}_{1}$ 与 $\mathcal{O}_{2}$ 为 $f(X)$ 的分离. 则 $f^{-1}(\mathcal{O}_{1})$ 与 $f^{-1}(\mathcal{O}_{2})$ 是 $X$ 中的不交非空开集, 其并是 $X$. 因此 $X$ 的这对分离与 $X$ 的连通性矛盾.
\end{proof}

我们把证明对实数集 $C$, 以下陈述
\begin{align}
\text{(i) }C\text{ 是区间;}\quad\text{(ii) }C\text{ 是凸的;}\quad\text{(iii) }C\text{ 是连通的}\label{align:ch11-连通-1}
\end{align}
等价留作练习.

\begin{definition}[介值性质]\label{definition:ch11-介值性质}
拓扑空间 $X$ 称为具有介值性质, 若 $X$ 上任何实值连续函数的象是一个区间.
\end{definition}

\begin{proposition}\label{proposition:ch11-连通介值}
拓扑空间具有介值性质当且仅当它是连通的.
\end{proposition}

\begin{proof}
根据\eqref{align:ch11-连通-1}, 实数的连通集是一个区间. 我们从命题 22 推出一个连通拓扑空间具有介值性质. 为证明逆命题, 我们假定 $X$ 是不连通的拓扑空间并得出它不具有介值性质. 事实上, 由于 $X$ 是不连通的, 存在一对 $X$ 的非空开子集 $\mathcal{O}_{1}$ 和 $\mathcal{O}_{2}$ 使得 $X=\mathcal{O}_{1}\cup\mathcal{O}_{2}$. 定义 $X$ 上的函数 $f$ 为在 $\mathcal{O}_{1}$ 上取值 $0$ 而在 $\mathcal{O}_{2}$ 上取值 $1$. 则 $f$ 是连续的, 由于对 $\mathbf{R}$ 的每个子集 $A$, 特别地, 对 $\mathbf{R}$ 的每个开子集 $A$, $f^{-1}(A)$ 是 $X$ 的开子集. 另一方面, $f$ 不具有介值性质.
\end{proof}

若一个拓扑空间是不连通的, 则对该空间的任何分离, 与这个分离中的每个集合有非空交的子空间也是不连通的. 此外, 实数区间在连续映射下的象是连通的. 因此拓扑空间 $X$ 是连通的, 若对每对点 $u,v\in X$, 存在连续映射 $f:[0,1]\to X$ 使得 $f(0)=u$ 且 $f(1)=v$. 具有该性质的拓扑空间称为弧连通的. 而弧连通拓扑空间是连通的, 存在不是弧连通的连通空间(见习题 49). 然而, 对欧氏空间 $\mathbf{R}^{n}$ 的开子集, 连通性等价于弧连通性(见习题 50).

\begin{exercise}\label{exercise:ch11-47}
令 $\{C_{\lambda}\}_{\lambda\in\Lambda}$ 为拓扑空间 $X$ 的连通子集族且假定它们中的任何两个有公共点. 证明 $\{C_{\lambda}\}_{\lambda\in\Lambda}$ 的并也是连通的.
\end{exercise}

\begin{exercise}\label{exercise:ch11-48}
令 $A$ 为拓扑空间 $X$ 的连通子集, 假定 $A\subseteq B\subseteq\overline{A}$. 证明 $B$ 是连通的.
\end{exercise}

\begin{exercise}\label{exercise:ch11-49}
证明平面的以下子集是连通的但不是弧连通的.
\begin{align*}
X=\{(x,y)\mid x=0,-1\leqslant y\leqslant 1\}\cup\{(x,y)\mid y=\sin 1/x,0<x\leqslant 1\}.
\end{align*}
\end{exercise}

\begin{exercise}\label{exercise:ch11-50}
证明弧连通拓扑空间 $X$ 是连通的. 也证明欧氏空间 $\mathbf{R}^{n}$ 的每个连通的开子集 $\mathcal{O}$ 是弧连通的. (提示: 令 $x$ 属于 $\mathcal{O}$. 定义 $C$ 为 $\mathcal{O}$ 中可在 $\mathcal{O}$ 内通过分段线性的弧连通 $x$ 的点集. 证明 $C$ 在 $\mathcal{O}$ 中既开又闭.)
\end{exercise}

\begin{exercise}\label{exercise:ch11-51}
考虑平面 $\mathbf{R}^{2}$ 上的圆 $C=\{(x,y)\mid x^{2}+y^{2}=1\}$. 证明 $C$ 是连通的.
\end{exercise}

\begin{exercise}\label{exercise:ch11-52}
证明 $\mathbf{R}^{n}$ 是连通的.
\end{exercise}

\begin{exercise}\label{exercise:ch11-53}
证明紧度量空间 $(X,\rho)$ 不是连通的, 当且仅当存在两个不交非空子集 $A$ 与 $B$, 其并是 $X$, 且存在 $\varepsilon>0$ 使得对所有 $u\in A$, $v\in B$, $\rho(u,v)\geqslant\varepsilon$. 证明对非紧度量空间这不成立.
\end{exercise}

\begin{exercise}\label{exercise:ch11-54}
度量空间 $(X,\rho)$ 称为是良链的, 若对每对点 $u,v\in X$ 与 $\varepsilon>0$, 存在 $X$ 中的有限个点 $u=x_{0},x_{1},\cdots,x_{n-1},x_{n}=v$ 使得对 $1\leqslant i\leqslant n$, $\rho(x_{i-1},x_{i})<\varepsilon$.
\begin{enumerate}[(i)]
\item 证明: 若 $X$ 是连通的, 则它是良链的, 但逆命题不成立.
\item 证明: 若 $X$ 是紧与良链的, 则它是连通的.
\item 证明: 若 $\mathbf{R}^{n}$ 的开子集是良链的, 则它是连通的.
\end{enumerate}
\end{exercise}

\begin{exercise}\label{exercise:ch11-55}
证明对平面 $\mathbf{R}^{2}$ 的任何点 $(x,y)$, 子空间 $\mathbf{R}^{2}\sim\{(x,y)\}$ 是连通的. 用此证明 $\mathbf{R}$ 不同胚于 $\mathbf{R}^{2}$.
\end{exercise}

\begin{exercise}\label{exercise:ch11-56}
证明\eqref{align:ch11-连通-1}中的三个断言的等价性.
\end{exercise}

\end{document}